\section{A four-level transport certificate}

\subsection{Exact and per-input certificates}

The first level searches a library of proved parameter transports. Structural
identity is the trivial member. Any nontrivial mapping must preserve the complete
transition, including reset and queues, and survive symbolic derivation,
exhaustive small-state checking, and adversarial trace search.

For a general bounded family, we unroll the recurrent network for $T$ steps and
propagate lower and upper state sets. Affine maps use sign-split interval
arithmetic. Fixed-point rounding is enclosed at integer bin boundaries; wrapping
that may cross the representable range expands to the full numeric domain.
Delays are queues of abstract states. At a threshold discontinuity, the analyzer
splits quiet and spiking branches, applies the appropriate reset, and conservatively
merges them. Discrete semantic variants are analyzed separately and their output
sets are joined.

For reference prediction $c_0$ and target logit intervals
$[\ell_c(x),u_c(x)]$, input $x$ is certified if
\begin{equation}
  \ell_{c_0}(x) > \max_{c\ne c_0} u_c(x).
\end{equation}
The condition proves an unchanged argmax for every semantics member represented
by the abstract family. Exact enumeration is also available when the declared
family itself is finite; a MILP checker is planned as a tighter small-network
oracle.

\subsection{Distribution certificate}

Static certification need not cover every input. On an untouched unlabeled audit
sample $x_1,\ldots,x_n$, define
\begin{equation}
  D_i(s)=\mathbf{1}[F_{s_0}(x_i)\ne F_s(x_i)].
\end{equation}
For $k_s=\sum_i D_i(s)$, we use the exact one-sided Clopper--Pearson limit
$U_s(\alpha)$. With $m$ target family members, Bonferroni correction at
$\alpha=\delta/m$ yields simultaneous coverage, and
$U_{\TRef,\Emu}=\max_s U_s(\delta/m)$ for a fixed but unknown declared target.

\begin{proposition}[Accuracy-change bound]
For arbitrary labels $Y$ and paired predictions from classifiers $f$ and $g$,
\[
|\Pr[f(X)=Y]-\Pr[g(X)=Y]|\leq \Pr[f(X)\ne g(X)].
\]
\end{proposition}
\begin{proof}
Pointwise, the two correctness indicators can differ only when the predictions
differ. Taking expectations and applying $|\mathbb{E}Z|\leq\mathbb{E}|Z|$
proves the statement.
\end{proof}

\subsection{Physical certificate}

An emulator cannot certify unmodeled hardware behavior. We therefore execute an
unlabeled canary and estimate paired emulator--hardware disagreement. Splitting
failure probability across the two estimates gives
\begin{equation}
U_{\mathrm{total}}=\min\{1,
U_{\TRef,\Emu}(\delta/2)+U_{\Emu,\Hw}(\delta/2)\}.
\end{equation}
Without the second term, the report is marked \emph{conditional on emulator}.
For nondeterministic hardware, the primary sample contains independent
$(x,\text{run})$ pairs. Repeated execution of selected inputs characterizes
run-to-run variance but cannot replace those pairs.
